\begin{abstract}
Verilog-A is the standard language for compact device models, yet the open tools that compile
it either target native code through LLVM or are tied to a host simulator. We present
\vas{}, a circuit simulator written from scratch in pure Rust that compiles Verilog-A source
into automatically differentiated model instances and solves them with a modified-nodal-analysis
Newton core: DC operating point and sweep, transient, AC and noise analyses. Two frozen internal
interfaces, an arena-based intermediate representation and a model-instance ABI, separate the
language front end from the numerical core. As a result the core is validated against
hand-written reference models, independently of the compiler.
Every analysis is checked against QSPICE golden outputs on 28 circuits to stated tolerances
($10^{-4}$ DC and AC, $10^{-3}$ transient and noise), and every differentiated operator against
finite differences. On the transistor-level ISCAS'85 suite built from PSP103 devices, all eleven
circuits, up to \num{268676} unknowns and \num{15882} PSP103 instances, reach operating
points whose every gate output matches the Boolean logic, and a c432 vector-sequence transient
matches on \num{1196800} settled samples. The simulator produces bit-identical output on Linux,
macOS and Windows and at any thread count. We describe the design choices that make this
determinism possible, the DC rescue strategy that makes the large circuits converge, and the
limitations that remain.
\todo{250-word limit; re-count after the numbers settle.}
\end{abstract}

\begin{IEEEkeywords}
Verilog-A, compact modeling, circuit simulation, automatic differentiation, modified nodal
analysis, reproducibility, Rust.
\end{IEEEkeywords}
